\chapter{Control flow and the execution mask}\label{ch:control-flow-execution}

A G17 simdgroup executes 32 lanes in lockstep. Divergence is handled by an execution state per lane:
instructions that change it (push a condition, flip it for an else, pop it), and two kinds of branch that move the program
counter only when the state of the whole simdgroup allows. Most conditionals do not branch; Apple's compiler
predicates them.

The execution state behaves as a per-lane nesting counter, active when zero, not as a
stack of masks (measured for the push and one-level pop; for the other forms, a model that balances every corpus program;
\cref{sec:execution-mask-family}). Structured control flow is built on that counter, and three of its consequences
first appeared as faulty programs. Code after a join runs on no lanes unless the join restores the
state (\cref{sec:execution-mask-family}). A loop nested inside a guarded region needs its own push before the loop, or the latch's pop removes the
region's level (\cref{sec:branches-back-edges}). A branch decides only whether the simdgroup skips work that no lane needs; which lanes run is set
by the counter (\cref{sec:predication-if-conversion,sec:branches-back-edges}).

Names used below:

\begin{itemize}
\item The exec family (one 4-byte encoding word, \cref{sec:execution-mask-family}): \op{582}; \op{577};
  \op{579}; the inverting if (\op{583}), the else forms (\op{575}, \op{576}) and the plain while
  (\op{578}), which have no glossary entries; \op{684}.
\item Branches (10 bytes): \op{458}, \op{462}, \op{450}; \op{580} ends a subroutine.
\item A \textbf{flag} is a per-lane predicate register (FLAG0-FLAG14) written by a compare such as \op{10369} (\cref{sec:compares-flags-selects}). FLAGTRUE is a constant source that always holds.
\item "Counter n" is the lane's nesting counter.
\end{itemize}

\section{Execution-mask instructions}\label{sec:execution-mask-family}

The family shares one encoding word (decoder; \code{agxforge/g17/cf.py}, \code{isa/g17-exec-mask.txt}). Bit 4 of byte 0 separates end of program
(\opn{684}, \code{0e 00 00 00}) from the family; bit 4 of byte 2 separates exec forms from branches; kind = b0.5 + 2 × b0.6; and
b2.1 inverts the predicate (\cref{tab:exec-family-encoding}).

\begin{xltabular}{\linewidth}{@{}lXXX@{}}
\caption{Exec-family and branch opcodes by kind, inversion bit b2.1 and branch bit b2.4.}\label{tab:exec-family-encoding}\\
\toprule
kind & b2.1 = 0 & b2.1 = 1 & with b2.4 = 1 (branch) \\
\midrule
\endfirsthead
\tablecontinued{4}
\toprule
kind & b2.1 = 0 & b2.1 = 1 & with b2.4 = 1 (branch) \\
\midrule
\endhead
0 if & \op{582} & \op{583} (no glossary entry) & stays exec (\op{573} / \op{574}, no glossary entries) \\*
1 pop & \op{577} & \opn{577} (b2.1 dead) & \op{462} \\*
2 else & \op{575} (no glossary entry) & \op{576} (no glossary entry) & \op{458} \\*
3 while & \op{578} (no glossary entry) & \op{579} & \op{450} \\*
\bottomrule
\end{xltabular}

\begin{itemize}
\item The count is b0.7 (high) and b2.3 (low), and it is a code table, not an integer: 00 → 1, 01 → 0, 10 → 2, 11 → 3. An
  encoder that writes the count as a plain 2-bit number emits \code{pop 0} when it means \code{pop 1}. Apple never emits count 3,
  nor a single pop deeper than 2.
\item The predicate is b1{[}4:6{]}, an 8-entry lookup: {[}74, 75, 76, 77, 78, 79, 3, 2{]}, register ids for FLAG0-FLAG5 and then two
  entries in another register file, which Apple uses to push a level unconditionally (FLAGTRUE). Pop reads no predicate.
\item Auxiliary operand 0 (bits b1.0, b1.1, b2.5-b2.7, b3.5, b3.7) takes only a few witnessed values per kind; its
  meaning is \emph{open}, and hard-coding one value writes an instruction Apple never emitted 4.6 percent of the time.
\item An encoder built from this field map re-encodes all 14,195 control-flow instructions in Apple's corpus byte for byte. It
  reached that only after the auxiliary field was carried as a field; as a constant it scored 94.34 percent.
\end{itemize}

\textbf{Worked example: an exec push and its join} (decoded; the second executed). \code{1e 00 00 0e}: b0 = 0x1e has bit 4 set, b0.5 = b0.6 = 0 (kind if),
b2 = 0x00 (exec, not inverted), count code 00 = 1: \op{582}, one level on FLAG0. \code{3e 03 40 0e}: b0 = 0x3e gives kind 1,
b2 = 0x40 has bit 4 clear, count code 00 = 1: \op{577}, pop 1. The second word is the join; it
was identified after a prediction failed (measured, \code{isa/g17-scalar-isa.toml}). A forward branch authored
one instruction past the compiler's own target, predicted to store 114, stored 0. A kernel was then built whose post-join
code writes a second value that does not depend on the branch (\cref{tab:join-removal}).

\begin{longtable}{@{}lll@{}}
\caption{Stored values with the join present and replaced by filler, on the taken and not-taken paths.}\label{tab:join-removal}\\
\toprule
join & taken (n = 2) & not taken (n = 7) \\
\midrule
\endfirsthead
\tablecontinued{3}
\toprule
join & taken (n = 2) & not taken (n = 7) \\
\midrule
\endhead
present & (114, 31) & (1868193279, 96) \\
replaced by filler & (0, 0) & (1868193279, 96) \\
\bottomrule
\end{longtable}

Without the join the branch-independent value was also lost, on the taken path only. With the output pre-filled with
0xDEADBEEF the sentinel remained, so the store did not execute. Code after a join therefore retires with no active lanes
unless the join restores the state, and a branch must target the join.

The counter semantics, per lane and with n the instruction's count, are these (measured for the push and one-level pop forms; whole-program for the rest; MM~\mm{25.145.2},
\code{agxforge/g17/releasecheck.py}):

\begin{itemize}
\item \textbf{if n}: an active lane (counter 0) stays 0 if its condition holds, else becomes n; an inactive lane adds n.
\item \textbf{else n}: 0 becomes n; n becomes 0 if the condition holds; any other value is unchanged.
\item \textbf{pop n}: counter = max(0, counter - n).
\item \textbf{while n}: an active lane whose condition fails becomes n (and waits for the matching pop n); otherwise unchanged.
\item FLAGTRUE always holds; the inverting if, the b2.1 = 1 else (no glossary entries) and while-exec, set mask in place
  negate the condition. A condition is the lane's flag value, which stays
  the same until a compare rewrites it.
\end{itemize}

Two observations rule out a push/pop stack in favour of the counter. Apple's corpus contains 338 top-level loops of the shape \code{while 2 / back edge / pop 2} with no push; a stack would pop levels never pushed, while the counter model balances all
6,594 corpus programs. And Apple over-pops on early-exit paths (in one probe a branch leaves a region two levels deep and
lands on \code{pop 2} then \code{pop 1}), which only a saturating counter tolerates; that one is an inference from the correctness of
Apple's code (decoder, \code{isa/g17-exec-mask.txt}). This is the model Asahi/Mesa call r0l. For every executed receipt that uses
a push and one-level pops, the counter and a stack agree. The in-place while and the two-level pop are needed by the
imageblock kernels' K loop and are admitted only on whole-program evidence. Nesting depth 1 to 6 with a store per level
emits one \code{if 1} per level and pops that sum to the pushes (depth 6: pop 2, pop 2, pop 2); else and while are net-neutral.

\emph{Correction:} the 4-byte backward branch in \code{isa/g17-scalar-isa.toml} (\code{de 6f f3 1e}) is the first four bytes of the
10-byte branch below (\code{agxforge/g17/cf.py}; \cref{app:a}). Author branches from the 10-byte map.

In the branch encoding (decoder, verified on 65 of 65 branches) b2.0 is set, and a signed two's-complement displacement has its
sign at b9.6 and weights 2 through $2^{47}$ scattered over bytes 0–9 (no weight-1 bit, so displacements are even). The target is the
branch's own offset + displacement, not the end of the instruction.

\section{Predication and if-conversion}\label{sec:predication-if-conversion}

Apple's compiler predicates first (Apple's compiler emits; \code{isa/g17-exec-mask.txt}, \code{isa/g17-scalar-isa.toml}). A
single-sided \code{if} emits no control-flow instruction: the store becomes a predicated store form, or the region is
if-converted to a compare-select (\cref{sec:compares-flags-selects}). Two or three nested levels with one assignment each still if-convert.
The exec machinery appears once a region has two sides, side effects at more than one level, or a loop. Forcing a real
branch took a conditional block of eight statements or more; below eight the block is predicated and no branch is emitted.

Apple places \op{462}
directly after a region's \op{582}; it is taken when no lane of the simdgroup is active and lands
on the region's pop, so it never changes which lanes run. On a 512-threadgroup grid with a selector matching no arm, both skips were taken and the program wrote
nothing (measured, MM~\mm{25.141.2}).

This project's compiler lowers a
conditional as pure predication by default (measured, MM~\mm{25.139.9}, MM~\mm{25.141.2}), so every dispatch walks the masked-off arms: +2.4~µs per dispatch for two arms
and +38~µs per four-op chain in an uber kernel. With the skip branch emitted (opt-in), an uber arm costs what the plain
kernel costs (131.2 against 131.0~µs; 68.8 against 68.7).

For lanes that are masked off:

\begin{itemize}
\item A release inside a region acts only on the lanes that executed the releasing instruction: with the region gated on
  t < 16, a released value read 0 on lanes 0–15 only, and with the region gated on t < 0 nothing was released (measured,
  MM~\mm{25.95}), so a value released in one arm is dead only on the lanes that took that arm. \Cref{ch:register-file} holds the release rules.
\item Load waits are per consumer, so a use after a join still waits for a load whose first consumer was skipped (this
  compiler, MM~\mm{25.141.2}).
\item Cross-lane reads inside a partially active simdgroup are a separate hazard (\cref{ch:cross-lane-operations}).
\end{itemize}

\section{Branches and loops}\label{sec:branches-back-edges}

The back edge, \op{458}, is taken iff any lane of the simdgroup is active. Every backward
branch in Apple's corpus is this branch (581 across 401 programs), every one of those programs also contains a while form, and
Apple never emits \op{450} backward (decoder, \code{isa/g17-exec-mask.txt}). \op{582} gates the
back edge; \op{579} sets the state in place and does not gate it (glossary). The loop shape Apple uses and this compiler emits is a latch \code{pop, compare, push, back edge}, with one more pop at the exit (measured, MM~\mm{25.139.9}).

A loop inside a region needs its own push, because the first trip's latch pop removes one level. At top level that is
harmless, but inside a guarded region it removes the region's level, and the guarded-off lanes then run everything after
the loop. The defect appeared when an uber kernel's masked-off arm stored threadgroup 0's rows; with a selector that matched no
arm, both arms stored them. The design measured for loops starts with a preheader push, which the lowering had never
emitted. This compiler now emits \code{0 < 1} and an exec push before a loop nested in a region, keeping every top-level loop's bytes
(measured, MM~\mm{25.139.9}).

A loop whose trip count is a runtime value lowers to an unsigned compare against a capped
count. The cap is required, since without it termination is not provable (measured, MM~\mm{25.90}). Apple's compiler emits \op{10369}/10 and \op{9328}/10 for such a loop
(MM~\mm{25.73}).

A loop-carried value must not be released. A release on the counter in the cap compare left a kernel unable to reach
either bound, and it ran until a reboot (measured, MM~\mm{25.116}; the case is described in
\cref{sec:loops-loop-variable-rule}). The same
family of defect, a loop phi's entry value released before the loop, made every thread store as thread 0 (MM~\mm{25.142.7},
MM~\mm{25.144.6}). The rules (a back edge reads its header phis' latch values; refuse a release of any loop-carried register
after the body's last write of it) are \cref{sec:loops-loop-variable-rule}'s, and the dispatch hazard is \cref{sec:forward-progress-in-order}'s.

\Cref{tab:loops-below-metal} lists loops executed below Metal (measured; programs authored by this compiler and submitted without Metal, each checked
against the exact expected output with guards, callbacks and recovery state; \code{docs/g17-first-dispatch-trials.md}).

\begin{xltabular}{\linewidth}{@{}>{\hsize=0.58\hsize}X>{\hsize=1.52\hsize}XX>{\hsize=0.84\hsize}X@{}}
\caption{Loop programs authored by this compiler and run below Metal, each checked against the exact expected output.}\label{tab:loops-below-metal}\\
\toprule
Trial & Program & Result & What it establishes \\
\midrule
\endfirsthead
\tablecontinued{4}
\toprule
Trial & Program & Result & What it establishes \\
\midrule
\endhead
Bounded per-lane GPU loop & 88 bytes: each lane starts at its lane index and adds 7 per trip; one-
and three-trip variants differ only at code byte 55 (compare immediate
0x80 against 0x82); \op{10369}, \op{582}, a 10-byte \op{458} at byte 62 with displacement −38 landing on
the body at byte 24, \op{577}, store, end & 64 of 64 words exact: 7..70 for one trip, 21..84 for three (the back
edge ran twice) & bounded cyclic control flow for this program and layout \\*
Per-lane divergent loop bounds & 146 bytes: bound loaded from A{[}lane{]}, runtime compare with a cap of
3, \op{458} at byte 120 with displacement
−82 to byte 38; byte-identical code under two input profiles & bounds 1,3 alternating gave 7, 22, 9, 24, \dots{} 69, 84; reversed gave 21,
8, 23, 10, \dots{} 83, 70: each lane computed lane + 7 × A{[}lane{]} beside
neighbours with different trip counts & per-lane divergent trip counts in one simdgroup \\*
Runtime loop cap and zero bound & the same 146 bytes & bounds 5,3: every lane i + 21 (over-cap lanes stopped at three trips);
bounds 0,3: even lanes i + 7, odd lanes i + 21 & the cap clamps; this emitted loop tests at the bottom, so a zero bound
still runs the body once \\*
Bounded loop inside a divergent region & 140 bytes: every lane stores its index, then only lanes with lane
< 32 enter a two-trip loop; pushes at offsets 50, 68 and 106,
pops at 96, 120 and 132, \op{458} at 110
with displacement −38 to the body at 72 & lanes 0–31 i + 14, lanes 32–63 i & the preheader push keeps the outer guard; not arbitrary nesting \\*
\bottomrule
\end{xltabular}

The zero-bound result comes from this loop shape and is not a machine property. A compiler that needs zero-trip
semantics must guard the loop or test before the body.

Apple's own program that uses call and \op{580} runs 32 of 32 lanes correctly in this
project's image path, and this compiler inlines calls (whole-program, MM~\mm{25.90}). Indirect calls are \emph{open}.

\leadhead{Safe assumptions.}

\begin{itemize}
\item The exec forms and branches Apple emits can be emitted with the field map of \cref{sec:execution-mask-family}; they
  are byte-identical to code that executed. No novel exec encoding (count 3, an unwitnessed auxiliary value, a predicate index Apple never selects) has been
  executed (\code{isa/g17-exec-mask.txt}).
\item The counter semantics hold for the exec push and the one-level exec pop (measured); for the in-place while and
  the two-level pop they hold only as whole programs use them. The else rule and the saturation are model statements consistent with all corpus programs, not
  separate measurements.
\item Branch-taken conditions (the skip branch when no lane is active, the back edge when any lane is) are measured for the
  skip-taken case and the executed loops; the skip branch's not-taken behaviour under partial activity follows from the
  same programs.
\end{itemize}

\textbf{Not known.} The auxiliary operand's meaning; where the counter is held and its width; indirect calls.

\textbf{Evidence.} MM~\mm{25.73}, MM~\mm{25.90}, MM~\mm{25.95}, MM~\mm{25.116}, MM~\mm{25.139.9}, MM~\mm{25.141.2}, MM~\mm{25.142.7}, MM~\mm{25.144.6}, MM~\mm{25.145.2};
\code{isa/g17-exec-mask.txt}, \code{isa/g17-scalar-isa.toml}, \code{agxforge/g17/cf.py}, \code{agxforge/g17/releasecheck.py},
\code{docs/g17-first-dispatch-trials.md}.
